The preceding parts of this thesis developed a systematic framework for
understanding the vacuum structure of quantum field theories and its
evolution with temperature. Beginning with the effective action
formalism in flat spacetime, we computed the one-loop effective
potential and saw how quantum corrections can qualitatively change the
vacuum structure of a theory, generating symmetry breaking where none
existed classically. Extending this to finite temperature, we found
that thermal corrections restore symmetry at high temperatures, driving
cosmological phase transitions as the early universe cooled. Throughout
all of this, the background spacetime was treated as fixed and flat, or
at most as a classical FRW geometry.

This assumption is an approximation, and in many physically important
situations it breaks down. In the very early universe, near the Planck
scale, the energy densities involved are large enough that the geometry
itself is dynamical and cannot be treated as a fixed background. Near
a collapsing star, the curvature of spacetime grows without bound as
the horizon forms. In the inflationary epoch, the exponential expansion
of the universe produces a de Sitter geometry whose causal structure is
fundamentally different from that of Minkowski space. In all of these
situations, the interplay between quantum matter and curved geometry
becomes essential.

The appropriate framework for addressing these questions is
\emph{semiclassical gravity}, in which the spacetime geometry is
treated classically via the Einstein equations, but the matter fields
are promoted to quantum operators. The backreaction of quantum matter
on the geometry is captured by the semiclassical Einstein equation,
\begin{equation}
    G_{\mu\nu} = 8\pi\langle\hat T_{\mu\nu}\rangle,
\end{equation}
where the right-hand side is the expectation value of the
quantum stress-energy tensor in some state. This is the natural
generalization of the classical field equations to include quantum
effects, and it is valid as long as the curvature remains well below
the Planck scale so that quantum gravitational effects can be ignored.

The necessity of this framework can be seen directly from the
effective action. In a curved background, the field equation acquires
a non-minimal coupling to the curvature,
$(\square - m^2 + \xi R)\phi = 0$, which shifts the effective mass of
the field and can qualitatively change the phase structure. More
fundamentally, the very notion of a vacuum state becomes
observer-dependent in curved spacetime. In flat spacetime, the
Poincar\'{e} invariance of the Minkowski metric singles out a unique
vacuum state agreed upon by all inertial observers, and the effective
potential is a well-defined quantity. In a general curved spacetime
there is no such preferred frame, and different observers, or the same
observer at different times, may disagree on the particle content of a
given state. The Bogoliubov transformation formalism, developed in
Chapter~\ref{ch:qft-curved}, makes this precise.

The physical consequences of this observer-dependence are profound.
The most celebrated is the Hawking effect: a black hole formed by
gravitational collapse emits a steady thermal flux of radiation at
a temperature $T_H = \kappa/2\pi$, where $\kappa$ is the surface
gravity of the horizon. This connects back to the cosmological context
in a direct way. The de Sitter spacetime produced during inflation also
possesses a horizon, the cosmological horizon, and quantum fields in
this background experience a Gibbons-Hawking temperature $T = H/2\pi$,
where $H$ is the Hubble rate. Quantum fluctuations in this thermal de
Sitter background seed the primordial density perturbations that
eventually grow into the large-scale structure of the universe we
observe today.

The chapters of this part develop the necessary tools in a logical
sequence. We begin with the causal structure of the Schwarzschild
spacetime, introducing the Kruskal and Penrose diagrams that make the
global geometry transparent. We then develop black hole thermodynamics,
establishing the Bekenstein-Hawking entropy and the generalized second
law. The general formalism of quantum field theory in curved spacetime
follows, including the Klein-Gordon inner product, the Bogoliubov
transformation, and the ambiguity in the notion of particles. We then
study two concrete examples of particle production, the moving mirror
and the uniformly accelerating detector, which isolate the essential
physics of the Hawking effect in a simpler setting. Finally we derive
Hawking radiation itself, both in the 2D model where the calculation
can be done exactly, and in the full 4D case where greybody factors
modify the spectrum.
